\documentclass[french]{yLectureNote}

\title{Fiche de révision}
\subtitle{Synthèse des notions essentielles}
\author{Paulhenry Saux}
\date{\today}
\yLanguage{Français}

\professor{}

\usepackage{graphicx}
\usepackage[utf8]{inputenc}
\usepackage{geometry}
\usepackage{tikz}
\usepackage{tkz-tab}
\usepackage{tabularx}
\usepackage{awesomebox}
\usepackage{menukeys}
\usepackage{fancyhdr}
\usepackage{blindtext}
\usepackage{hyperref}
\usepackage{caption}
\usepackage{pifont}
\usepackage{array}
\usepackage{lipsum}
\usepackage{yFlatTable}
\usepackage{multicol}

\newcommand{\Lim}[1]{\lim\limits_{\substack{#1}}\:}
\renewcommand{\vec}{\overrightarrow}
\newcommand{\bz}{\overline{z}}
\DeclareMathOperator{\card}{card}
\renewcommand{\vec}{\overrightarrow}
\newcommand{\N}[0]{\mathbb{N}}
\newcommand{\R}[0]{\mathbb{R}}
\newcommand{\C}[0]{\mathbb{C}}
\newcommand{\dd}[0]{\mathrm{d}}
\newcommand{\er}{\vec{e_{\rho}}}
\newcommand{\ep}{\vec{e_{\varphi}}}
\newcommand{\pip}{\pi^+}
\newcommand{\pim}{\pi^-}
\newcommand{\mv}{\mathcal{V}}
\DeclareMathOperator{\Div}{Div}
\newcommand{\rot}{\vec{\mathrm{rot}}}
\newcommand{\grad}{\vec{\mathrm{grad}}}

\begin{document}
\chapter{Produit scalaire}
\section{Produit scalaire et norme}

\begin{definition}
[Produit scalaire]
    Soit $E$ un ev sur $\R$ de dimebsion finie ou non.
    $$\varphi : E\times E \to \R, (x,y)\to \varphi(x,y)$$
    \begin{itemize}
        \item $\varphi$ est une forme (l'espace d'arrivée est $\R$) bilinéaire si $x\to \varphi(x,y)$ est linéaire, de m\^eme pour y
        \item $\varphi$ est symétriuque $\iff \varphi(x,y) = \varphi(y,x)$
        \item $\forall x, \varphi(x,x) \geq0$
        \item $\varphi$ est définie positive $\iff \forall x, \varphi(x,x)\geq 0 \wedge \varphi(x,x) = 0 \Rightarrow x=0$
    \end{itemize}
    On dit que $\varphi$ est un produit scalaire $\iff \varphi$ est une forme bilinéaire symétrique définie positive.
\end{definition}
\begin{definition}
[Espace Euclidien]
    SI la dimension de $E$ est finie et $\varphi$ produit scalaire, alors la paire $(E,\varphi)$ est un espace euclidien.


\end{definition}

\begin{definition}
[espace Préhilbertien]
        Si la dimension de $E$ est infinie, et $\varphi$ produit scalaire, alors la paire $(E,\varphi)$ est un espace préhilbertien.

\end{definition}
\begin{definition}
[Norme]
    On note la norme associée au produit scalaire $\sqrt{(x|x)}$.

\end{definition}

\begin{proposition}
[Inégalité de Cauchy Swartz]
    On munit E d'un produit scalaire, alors $\forall(x,y), (x|y) \leq ||x||||y||$

\end{proposition}
\marginCheck[Cas d'égalité]{ Le produit de deux vecteurs ne peut pas dépasser le produit de leurs longueurs. L'égalité n'arrive que si x et y sont colinéaires.}
\begin{proposition}
[Propriétés de la norme]
    \begin{itemize}
        \item $\forall x\in E, ||x||\geq 0$
        \item $||x|| = 0 \Rightarrow x=0$
        \item $||\lambda x|| = |\lambda|||x||$
        \item $||x+y|| \leq ||x||+||y||$
    \end{itemize}

\end{proposition}

\begin{proposition}
[Calcul du produit scalaire]
    \begin{itemize}
        \item $(x|y) = \frac12 (||x+y||^2-||x||^2-||y||^2)$
        \item $(x|y) = \frac12(-||x-y||^2+||x||^2+||y||^2)$
        \item $(x|y) = \frac14(||x+y||^2-||x-y||^2)$
        \item $||x+y||^2+||x-y||^2 = 2(||x||^2+||y||^2)$
    \end{itemize}

\end{proposition}

\section{Orthogonalité}

\begin{definition}
[Vecteurs orthogonaux]
    Deux vecteurs y et x sont orthogonaux quand leur produit scalaire est nul. On note $x\perp y$

\end{definition}

\begin{theorem}
[Pythagore]
    En réels, $x\perp y\iff ||x+y||^2 = ||x||^2+||y||^2$

\end{theorem}
\marginCritical[Validité]{Attention, Pythagore ne marche sous cette forme simple que dans les espaces réels.}
\begin{definition}
[Base orthogonale/orthonormée]
Une base $e_1,\dots,e_n$ est orthogonale quand $\forall i,j, i\neq j, e_i\perp e_j$ 

    Une base est orthonormée quand $\forall i\neq j, e_i\perp e_j, \forall i, ||e_i||=1$

\end{definition}

\begin{proposition}
[Décomposition dans une base orthonormée]
Si $e_1\dots e_n$ est un BON, alors $\forall x\in E, x= (e_1|x)e_1+\dots + (e_n|x)e_n$
    Les coordonnées de x dans la base $e_1,\dots, e_n$, notées $(x_1,\dots, x_n)$ sont données par $x_i = (e_i|x)$

\end{proposition}

\begin{proposition}
[Procédé de Gram-Schmidt]
Soit $(u_1, _dots, u_n)$ une base quelconque de $E$. Il existe une base orthonormée $e_1,\dots, e_n$
    telle que $\forall p, Vect(e_1,\dots, e_p) = Vect(u_1,\dots, u_p)$

\end{proposition}

\section{Sous espace orthogonal et projection}

\begin{definition}
[Orthogonal d'une partie]
Soit E un espace euclidien. A est une partie de E. L'orthogonal de A, noté 
    $A^{\perp}$ est donné par $A^{\perp} = \{y\in E/ \forall a \in A (y|a)=0\}$

\end{definition}

\begin{proposition}
[Orthogonal d'une partie : sous-espace]
$A^{\perp}$ est toujours un sev de E m\^eme si A n'est qu'une partie quelconque de E.

\end{proposition}

\begin{proposition}[Supplémentaire orthogonal]
[Supplémentaire orthogonal]
$F$ sev de E euclicien, F et $F^{\perp}$ sont supplémentaires $\iff E = F\oplus F^{\perp}$

\end{proposition}

\marginCritical[Validité]{En dimension infinie, cette propriété (Supplémentaire orthogonal) est fausse en général !}

\begin{proposition}
[Projection orthogonale]
F un SEV de E, $x\in E$. Alors $\forall y \in F, ||x-p_f(x)||\leq||x-y|| \iff ||x-p_f(x)|| = \min(||x-y||)$
    C'est la distance de x à F et $p_f(x)$ est la projection orthogonale de x sur F.

\end{proposition}

\end{document}
